\section{Related Work}
\label{sec:related}

\para{Representation change under fine-tuning.} \citet{merchant2020happens} showed that fine-tuning BERT mostly changes the upper layers and leaves lower layers close to the pretrained model. Representation similarity is usually measured globally, with SVCCA~\citep{raghu2017svcca} or centred kernel alignment~\citep{kornblith2019similarity}. Global measures can hide large changes in local structure, because a near-rotation of the space keeps CKA high while neighbours change. We therefore focus on kNN neighbourhood preservation, a local measure common in dimensionality-reduction evaluation~\citep{venna2006local}, and report CKA alongside it. Unlike earlier end-of-training comparisons, we follow the full trajectory over 17 checkpoints.

\para{Anisotropy of contextual embeddings.} Contextual representations occupy a narrow cone~\citep{ethayarajh2019contextual}, and removing dominant directions improves static embeddings~\citep{mu2018allbutthetop}. \citet{rudman2022isoscore} proposed IsoScore to measure how uniformly a space uses its dimensions. These works describe the isotropy of a trained model. We ask whether fine-tuning moves isotropy in a consistent direction, and find that it does not.

\para{Forgetting in language-model fine-tuning.} Catastrophic forgetting has been studied since early connectionist work~\citep{mccloskey1989catastrophic,kirkpatrick2017overcoming}. For LLMs, \citet{luo2023empirical} measured forgetting during continual instruction tuning, \citet{kalajdzievski2024scaling} fitted scaling laws relating forgetting to fine-tuning loss and parameter count, and \citet{biderman2024lora} showed that LoRA~\citep{hu2022lora} learns less and forgets less than full fine-tuning. Our learning-rate and full fine-tuning ablations agree that update size, not only scale, sets how much is lost. We add a link between local geometry and item-level forgetting, and show that the geometry-to-forgetting mapping differs across model families.

\para{Domain distance and transfer.} Domain-adaptive pretraining work~\citep{gururangan2020dont} uses vocabulary overlap to describe how far a domain is from pretraining data. We test whether this distance predicts how much general-text geometry fine-tuning preserves, and find the opposite of the intuitive prediction.
